% Root and B literal audits passed; grid endpoints explicitly repaired.
\subsection{Two minimal rows at each object of an act source}

\begin{lemma}[Finite affine diagrams on a common rational grid]
\label{C:categorical-grid-act-cups}
Let $D$ be a finite category and let $T:D^{\mathrm{op}}\to\mathrm{Set}$
be a finite right act. Suppose each $T_i$ is the set of all points
of $\delta\mathbb Z$ in a nonempty closed interval $I_i$ whose
endpoints lie in $\delta\mathbb Z$, where $\delta>0$ is fixed and rational. Suppose every act map is the
restriction of an affine map $I_j\to I_i$ of slope $0,1$, or $-1$.
Degenerate intervals are allowed. Then every positive-degree
rational cohomology product of the act cone is zero.
\end{lemma}
\begin{proof}
The affine maps give a diagram of contractible intervals.
For a nondegenerate source interval, their composition identities
follow from those on its two endpoints. For a degenerate source
interval there is only one point to check. Thus this is an actual
diagram, and its homotopy colimit maps equivalently to $BD$.
The pointed homotopy colimit of the objectwise quotients $I_i/T_i$
is therefore the cofiber of $\operatorname{hocolim}T\to BD$,
which is the classifying space of the act cone. As in
Lemma~\ref{C:two-states-objectwise-act-cups}, this identification
collapses the basepoint section of the ordinary category bar.

Each nondegenerate quotient is a wedge of circles indexed by
consecutive grid intervals of length $\delta$. An affine act map
of slope $1$ or $-1$ sends each such interval onto one interval
of the same length, preserving or reversing its orientation.
A slope-zero map becomes the pointed constant map. The
degree-two map on every circle, using its interval parameter
modulo endpoints, consequently commutes strictly with every act
map. It is also compatible with degenerate intervals.
The reduced cellular chain complex of each quotient is concentrated
in degree one, and this natural self-map is multiplication by two
on it. On the reduced pointed category bar it is multiplication
by two in every total degree. It thus acts by two on all positive
rational cohomology. Multiplicativity gives $2xy=4xy$, proving
$xy=0$ for any two positive-degree classes.
\end{proof}

\begin{theorem}[The categorical two-row act obstruction]
\label{C:categorical-two-row-act-cups}
Let $S$ be a finite strongly connected category. If every object
has at most two rows in its minimal Rees ideal, then every product
of rational degree-two classes on the cone of any nonempty finite
right $S$-act is zero. There is no restriction on the maximal
group of the minimal ideal or on the number of states. Dually,
at most two minimal columns at each object exclude such products
on every finite left-act cone.
\end{theorem}
\begin{proof}
Put $R=\mathbb QS$, and write $\mathbf1$ for the constant
left $S$-module. Choose a minimal idempotent at an object $a$,
let $G$ be its maximal group, and put
\[
 \bar e=\frac1{|G|}\sum_{g\in G}g,\qquad P=R\bar e.
\]
This is a projective left $R$-module. At an object $i$ its
basis consists of the $G$-averaged rows of $S(a,i)$ in the minimal
ideal. Every arrow sends these row basis vectors to row basis
vectors. Objectwise augmentation gives an exact sequence
\[
 0\longrightarrow V\longrightarrow P\longrightarrow\mathbf1
 \longrightarrow0.
\]
Here $V_i$ has dimension zero or one. Choose the usual difference
basis when there are two rows. For $f:i\to j$, its action on
the difference spaces is a scalar $\chi_f\in\{0,1,-1\}$,
with zero used whenever either space is zero. These are typed
module maps and satisfy the corresponding composition identities.

Let $F:S^{\mathrm{op}}\to\mathrm{Set}$ be the right act.
All its stalks are nonempty: strong connectivity propagates
nonemptiness along the contravariant maps. Put
$L_F=\ker(\mathbb QF\to\mathbf1_{\mathrm{right}})$.
The augmented category bar, followed by its cofiber, gives
\[
 H^2(B\mathcal C(S,F);\mathbb Q)
   \cong D\operatorname{Tor}^{R}_1(L_F,\mathbf1)
   \cong\operatorname{Ext}^{1}_R(\mathbf1,DL_F).
\]
These are the objectwise versions of
Lemma~\ref{C:act-cone-Tor}; explicitly the cofiber chain complex
is the suspension of $L_F\otimes_R^{\mathbf L}\mathbf1$.
Finite-dimensional vector-space duality gives the second equality.
Projectivity of $P$ gives a natural surjection
\[
 \operatorname{Hom}_R(V,DL_F)\longrightarrow
 H^2(B\mathcal C(S,F);\mathbb Q).
\]

Choose a lift of any class $x$ through this surjection. It is a
family of functionals $\ell_i:L_F(i)\to\mathbb Q$, with
$\ell_i=0$ if $V_i=0$, satisfying
\[
 \ell_i L_F(f)=\chi_f\ell_j\qquad(f:i\to j).
\]
Choose $s_i\in F(i)$ and set
$q_i(s)=\ell_i([s]-[s_i])$. There are rational numbers $\beta_f$
such that
\[
 q_i(F(f)s)=\chi_fq_j(s)+\beta_f.
\]
The equation follows by subtracting the value at $s_j$; one may
take $\beta_f=q_i(F(f)s_j)$. At a zero-dimensional $V_i$ the
set of values is the singleton $\{0\}$, which causes no exception.

Choose a positive rational $\delta$ for which all the finitely
many values $q_i(s)$ and $\beta_f$ lie in $\delta\mathbb Z$.
Let $I_i$ be the convex hull of $q_i(F(i))$, and enlarge that
finite set to $T_i=I_i\cap\delta\mathbb Z$. Each affine map
$t\mapsto\chi_ft+\beta_f$ carries $I_j$ into $I_i$ and the
whole grid $T_j$ into $T_i$. The typed composition identities
hold on these intervals: if $I_j$ is nondegenerate, equality
at its endpoints determines the affine map, and if it is
degenerate there is just one value to check. Thus the $T_i$
form a finite right $S$-act and $q:F\to T$ is an actual act map.

On $L_T(i)$ the coordinate functional
\[
 \ell_i^T\left(\sum_{t\in T_i}a_t[t]\right)=\sum_{t\in T_i}a_t t,
 \qquad \sum_{t\in T_i}a_t=0,
\]
is zero when $V_i=0$ and otherwise satisfies the same typed
equivariance equation. The translation disappears because the
coefficients sum to zero. Moreover $\ell_i=\ell_i^T L(q_i)$.
The natural connecting map therefore defines a class $y$ on
$B\mathcal C(S,T)$ with $x=(B\mathcal C(S,q))^*y$.
Lemma~\ref{C:categorical-grid-act-cups} gives $y^2=0$, hence
$x^2=0$. This applies to every $x$. Rational polarization proves
that every mixed degree-two product also vanishes. Taking
opposites proves the assertion about columns and left acts.
\end{proof}

\begin{corollary}\label{C:categorical-two-row-projective-sink}
The same degree-two product obstruction applies when the trivial
sink of the preceding theorem is replaced by a finite monoid $N$
whose trivial right rational module is projective, and the cross
sets form compatible $(N,S)$-actions. No finite-type hypothesis
is required for this assertion.
\end{corollary}
\begin{proof}
Proposition~\ref{C:finite-target-extreme-corner} gives an actual
quotient functor replacing $N$ by a point and every incident
left $N$-set by its orbit set. It is a rational cohomology-ring
equivalence. Apply Theorem~\ref{C:categorical-two-row-act-cups}
to the resulting right $S$-act.
\end{proof}

\begin{proposition}\label{C:strong-pair-fourteen-source-cone}
Let $S$ be a finite strongly connected category on two objects
with at most fourteen arrows. Adjoin a trivial sink with a
nonempty finite right $S$-act. If the resulting category has
at most twenty-one arrows, it cannot have type $\mathbb CP^2$.
The dual assertion holds for an adjoined trivial source.
\end{proposition}
\begin{proof}
If one object of $S$ is a retract of the other, deleting it
preserves the whole category's homotopy type and gives an
ordinary act cone with fewer than twenty-two arrows. This is
excluded by Corollary~\ref{C:act-cone-minimum-twenty-two}.
Assume henceforth that neither object is such a retract.
Both corner identities lie outside the minimal ideal of $S$.

If $H^2(BS;\mathbb Q)=0$, every degree-two class of the act cone
is relative to $BS$. Its relative quotient is a suspension, so
all degree-two products vanish. Otherwise the categorical Rees
comparison requires both total minimal row and column counts,
say $r,c$, to be at least three. If the minimal group had order
at least two, its ideal and the two outside identities would
already require at least $2\cdot3\cdot3+2=20$ arrows.
The minimal group is therefore trivial, and $rc+2\leq14$.
The only possibilities are $(r,c)=(3,3),(3,4),(4,3)$.

If each object has at most two rows, apply
Theorem~\ref{C:categorical-two-row-act-cups}.
The only remaining case has row counts $(3,1)$, total column
count three, and exactly fourteen arrows. Every arrow other
than the two identities is in the minimal band. Let $b$ be
the one-row object and choose a minimal idempotent $e_b$ there.
Every incoming arrow from the other object lies in the minimal
band and is fixed by left multiplication by $e_b$. The corner
at $b$ has no unit other than its identity. Minimal cross arrows
in the two directions have composite $e_b$.
There is no external entrance at $b$, because the added object
is a sink. Lemma~\ref{C:one-sided-corner-deletion} consequently
deletes $b$ from the whole category by a homotopy equivalence.

The remaining source monoid consists of the identity and a
$3\times c_a$ band, where $c_a=1$ or $2$. If $c_a=1$, its
classifying space is contractible because the band elements
are left zeros, so the preceding relative-suspension argument
applies. If $c_a=2$, it has only one element outside its
two-column band, and
Lemma~\ref{C:three-outside-all-two-column-act-cups} makes all
degree-two products of every act cone zero. Both alternatives
contradict the nonzero projective-plane square. Opposition
proves the dual assertion.
\end{proof}


\begin{corollary}\label{C:two-valued-representable-act-cups}
Let $S$ be a finite category with an object $a$ for which
$1\leq|S(a,i)|\leq2$ at every object $i$. For every finite right
$S$-act whose stalks are all nonempty, every rational degree-two
product of its act cone is zero.
\end{corollary}
\begin{proof}
Repeat the proof of Theorem~\ref{C:categorical-two-row-act-cups}
with the projective representable module $P=\mathbb QS(a,-)$.
Its row basis at each object has one or two elements, and
composition sends basis elements to basis elements. Thus its
augmentation kernel has the same objectwise zero- or
one-dimensional form. The Ext lift, finite rational grid, actual
act map, and scalar-degree argument apply verbatim. Strong
connectivity was used there only to construct this projective
cover and to ensure nonempty stalks, which are now explicit
hypotheses.
\end{proof}

\begin{corollary}\label{C:two-orbit-projective-cover-act-cups}
Let $S$ be a finite category, $a$ an object, $e\in\operatorname{End}_S(a)$
an idempotent, and $G$ a finite subgroup of the unit group of
$e\operatorname{End}_S(a)e$. Suppose that the right $G$-orbit
set $S(a,i)e/G$ has one or two elements at every object $i$.
For every finite right $S$-act with all stalks nonempty, all
rational degree-two products of its act cone vanish.
The conclusion remains true after replacing the trivial sink by
a monoid with projective trivial right rational module.
\end{corollary}
\begin{proof}
The averaging idempotent $\bar e=|G|^{-1}\sum_{g\in G}g$ gives
the projective module $P=\mathbb QS\bar e$. At $i$, a basis
consists of the normalized averages of the right $G$-orbits
in $S(a,i)e$. The action need not be free. Left composition
maps one orbit equivariantly onto another, and sends its normalized
average to the normalized average of that target orbit.
Consequently the augmentation $P\to\mathbf1$ is precisely
the projective basis cover used in the proof of
Theorem~\ref{C:categorical-two-row-act-cups}, with one or two
basis elements at each object. That proof applies unchanged.
For the last assertion, the actual extreme-corner orbit functor
is a rational cohomology-ring equivalence and does not increase
any cross stalk. Apply the first assertion to its act cone.
\end{proof}

